\documentclass[12pt]{article}
\makeatletter
\def\input@path{{../common/}{../guide/}{../data/}{../code/}}
\makeatother
\input{preamble}
\pagestyle{fancy}
\setlength{\headheight}{42pt}

 
\begin{document}
\lhead{Ingeniería Física \\ Escuela de Física \\ Tecnológico de Costa Rica} 
\rhead{Instrumentación I \\ Tarea \#1  \\ Entrega: Semana 3} 
\cfoot{\thepage\ de \pageref{LastPage}}
\setlength{\parindent}{0em}

1. A partir del diagrama incluido en la Figura \ref{t1f1} y usando el estándar \href{https://www.isa.org/products/ansi-isa-5-1-2024-instrumentation-and-control-symb}{ANSI/ISA-5.1-2024}.
Realice los siguiente (50 pts):
\begin{enumerate}
    \item Para cada uno de los elementos de la figura que estan incluidos en el estándar incluya*: 
    \begin{enumerate}
        \item Símbolo
        \item Breve descripción
        \item Variable a medir o controlar en caso de instrumentos de medición o control
        \item Tabla de la norma donde aparece. Por ejemplo: Tabla 5.2.4        
    \end{enumerate}
\end{enumerate}
\begin{figure}[H]
    \centering
    \includegraphics[width = 0.6\linewidth]{fig/diagrama1.png}
    \caption{Separador de tres fases}
    \label{t1f1}
\end{figure}
* Si algún elemento aparece varias veces se debe incluir una sola vez

2. Investigue el funcionamiento de las siguientes funciones en Arduino (30 pts):

\begin{enumerate}
    \item pinMode()
    \item digitalRead()
    \item digitalWrite()
    \item analogRead()
    \item analogWrite()
    \item delay()
    \item Serial.begin()
    \item Serial.print()
    \item Serial.readString()
    \item Serial.available()
\end{enumerate}

3. Escoja un sistema físico que usted conozca (péndulo simple, un circuito eléctrico o un sistema hidráulico por ejemplo), explíquelo y haga el modelo matemático del mismo (20 pts).

\end{document}